\section{账本事件的区域史深描：western-xia}

\label{sec:ledger-depth-western-xia}

本组把账本事件置于战争、制度、区域社会与证据类型的交叉处。年表提供政治节点，地方材料与物质遗存则限制我们对征发、身份、信仰和迁徙所能作出的断言。

\subsection{战争、行政与区域资源}

账本条目\texttt{XIA-0982-EVENT-049}所记982年的“李继迁脱离宋朝羁縻控制，在西北边地集结力量”，属于党项李氏与宋的历史语境。条目所示前因是前期边防、财政和统治联盟的压力构成直接背景，具体条件须按本条材料分辨，其可见后果为改变后续资源配置与政治选择；实际执行范围和社会影响仍须分项核验。这一变化若进入区域生活，必须通过道路、文书、军队、市场、家庭和地方中介才会发生；政权易手或法令发布不等于所有城镇、乡村与边地在同一时间获得同样经验。

本条所列来源为《宋史》卷485〈夏国传〉；《辽史》《金史》相关本纪；黑水城西夏文文书与考古报告。它们能够支持特定的年序、制度语言、政治行动或局部遗存，却不能无条件化为社会总量。账本特别提醒“编年主干可靠；细节、规模与动机须与对方材料、地方文书或考古材料互校”。因此，军报的兵数、条约的名分、官府的族群分类、判牍中的道德判断以及后出的叙述都须回到产生它们的场景，而不可直接代表全域动机或人口。

将事件与地方层次连接时，应同时考察资源怎样被征集、信息怎样传递、迁徙怎样改变户籍和劳动、宗教与亲属网络怎样提供或限制协作。现有证据往往只照亮少数地点；保留这种不均衡，才能避免把宋辽夏金元的多政权历史改写成单一中心的必然过程。

账本条目\texttt{XIA-0982-REGNAL-YEAR}所记982年的“党项李氏以李继迁在位第1年作为诏令、赋役与外交文书的纪年框架”，属于党项李氏的历史语境。条目所示前因是新岁颁历、年号和纪年是中枢与地方行政沟通的共同前提，其可见后果为为同年政令、战事和地方材料提供日期锚点，不能单独推知具体政策执行。这一变化若进入区域生活，必须通过道路、文书、军队、市场、家庭和地方中介才会发生；政权易手或法令发布不等于所有城镇、乡村与边地在同一时间获得同样经验。

本条所列来源为《宋史》卷485〈夏国传〉；《辽史》《金史》相关本纪；黑水城西夏文文书与考古报告。它们能够支持特定的年序、制度语言、政治行动或局部遗存，却不能无条件化为社会总量。账本特别提醒“纪年框架可靠；该记录仅确证政权纪年连续，年内具体事项须由本年条另证”。因此，军报的兵数、条约的名分、官府的族群分类、判牍中的道德判断以及后出的叙述都须回到产生它们的场景，而不可直接代表全域动机或人口。

将事件与地方层次连接时，应同时考察资源怎样被征集、信息怎样传递、迁徙怎样改变户籍和劳动、宗教与亲属网络怎样提供或限制协作。现有证据往往只照亮少数地点；保留这种不均衡，才能避免把宋辽夏金元的多政权历史改写成单一中心的必然过程。

账本条目\texttt{XIA-0983-REGNAL-YEAR}所记983年的“党项李氏以李继迁在位第2年作为诏令、赋役与外交文书的纪年框架”，属于党项李氏的历史语境。条目所示前因是新岁颁历、年号和纪年是中枢与地方行政沟通的共同前提，其可见后果为为同年政令、战事和地方材料提供日期锚点，不能单独推知具体政策执行。这一变化若进入区域生活，必须通过道路、文书、军队、市场、家庭和地方中介才会发生；政权易手或法令发布不等于所有城镇、乡村与边地在同一时间获得同样经验。

本条所列来源为《宋史》卷485〈夏国传〉；《辽史》《金史》相关本纪；黑水城西夏文文书与考古报告。它们能够支持特定的年序、制度语言、政治行动或局部遗存，却不能无条件化为社会总量。账本特别提醒“纪年框架可靠；该记录仅确证政权纪年连续，年内具体事项须由本年条另证”。因此，军报的兵数、条约的名分、官府的族群分类、判牍中的道德判断以及后出的叙述都须回到产生它们的场景，而不可直接代表全域动机或人口。

将事件与地方层次连接时，应同时考察资源怎样被征集、信息怎样传递、迁徙怎样改变户籍和劳动、宗教与亲属网络怎样提供或限制协作。现有证据往往只照亮少数地点；保留这种不均衡，才能避免把宋辽夏金元的多政权历史改写成单一中心的必然过程。

账本条目\texttt{XIA-0984-REGNAL-YEAR}所记984年的“党项李氏以李继迁在位第3年作为诏令、赋役与外交文书的纪年框架”，属于党项李氏的历史语境。条目所示前因是新岁颁历、年号和纪年是中枢与地方行政沟通的共同前提，其可见后果为为同年政令、战事和地方材料提供日期锚点，不能单独推知具体政策执行。这一变化若进入区域生活，必须通过道路、文书、军队、市场、家庭和地方中介才会发生；政权易手或法令发布不等于所有城镇、乡村与边地在同一时间获得同样经验。

本条所列来源为《宋史》卷485〈夏国传〉；《辽史》《金史》相关本纪；黑水城西夏文文书与考古报告。它们能够支持特定的年序、制度语言、政治行动或局部遗存，却不能无条件化为社会总量。账本特别提醒“纪年框架可靠；该记录仅确证政权纪年连续，年内具体事项须由本年条另证”。因此，军报的兵数、条约的名分、官府的族群分类、判牍中的道德判断以及后出的叙述都须回到产生它们的场景，而不可直接代表全域动机或人口。

将事件与地方层次连接时，应同时考察资源怎样被征集、信息怎样传递、迁徙怎样改变户籍和劳动、宗教与亲属网络怎样提供或限制协作。现有证据往往只照亮少数地点；保留这种不均衡，才能避免把宋辽夏金元的多政权历史改写成单一中心的必然过程。

账本条目\texttt{XIA-0985-REGNAL-YEAR}所记985年的“党项李氏以李继迁在位第4年作为诏令、赋役与外交文书的纪年框架”，属于党项李氏的历史语境。条目所示前因是新岁颁历、年号和纪年是中枢与地方行政沟通的共同前提，其可见后果为为同年政令、战事和地方材料提供日期锚点，不能单独推知具体政策执行。这一变化若进入区域生活，必须通过道路、文书、军队、市场、家庭和地方中介才会发生；政权易手或法令发布不等于所有城镇、乡村与边地在同一时间获得同样经验。

本条所列来源为《宋史》卷485〈夏国传〉；《辽史》《金史》相关本纪；黑水城西夏文文书与考古报告。它们能够支持特定的年序、制度语言、政治行动或局部遗存，却不能无条件化为社会总量。账本特别提醒“纪年框架可靠；该记录仅确证政权纪年连续，年内具体事项须由本年条另证”。因此，军报的兵数、条约的名分、官府的族群分类、判牍中的道德判断以及后出的叙述都须回到产生它们的场景，而不可直接代表全域动机或人口。

将事件与地方层次连接时，应同时考察资源怎样被征集、信息怎样传递、迁徙怎样改变户籍和劳动、宗教与亲属网络怎样提供或限制协作。现有证据往往只照亮少数地点；保留这种不均衡，才能避免把宋辽夏金元的多政权历史改写成单一中心的必然过程。

账本条目\texttt{XIA-0986-REGNAL-YEAR}所记986年的“党项李氏以李继迁在位第5年作为诏令、赋役与外交文书的纪年框架”，属于党项李氏的历史语境。条目所示前因是新岁颁历、年号和纪年是中枢与地方行政沟通的共同前提，其可见后果为为同年政令、战事和地方材料提供日期锚点，不能单独推知具体政策执行。这一变化若进入区域生活，必须通过道路、文书、军队、市场、家庭和地方中介才会发生；政权易手或法令发布不等于所有城镇、乡村与边地在同一时间获得同样经验。

本条所列来源为《宋史》卷485〈夏国传〉；《辽史》《金史》相关本纪；黑水城西夏文文书与考古报告。它们能够支持特定的年序、制度语言、政治行动或局部遗存，却不能无条件化为社会总量。账本特别提醒“纪年框架可靠；该记录仅确证政权纪年连续，年内具体事项须由本年条另证”。因此，军报的兵数、条约的名分、官府的族群分类、判牍中的道德判断以及后出的叙述都须回到产生它们的场景，而不可直接代表全域动机或人口。

将事件与地方层次连接时，应同时考察资源怎样被征集、信息怎样传递、迁徙怎样改变户籍和劳动、宗教与亲属网络怎样提供或限制协作。现有证据往往只照亮少数地点；保留这种不均衡，才能避免把宋辽夏金元的多政权历史改写成单一中心的必然过程。

账本条目\texttt{XIA-0987-REGNAL-YEAR}所记987年的“党项李氏以李继迁在位第6年作为诏令、赋役与外交文书的纪年框架”，属于党项李氏的历史语境。条目所示前因是新岁颁历、年号和纪年是中枢与地方行政沟通的共同前提，其可见后果为为同年政令、战事和地方材料提供日期锚点，不能单独推知具体政策执行。这一变化若进入区域生活，必须通过道路、文书、军队、市场、家庭和地方中介才会发生；政权易手或法令发布不等于所有城镇、乡村与边地在同一时间获得同样经验。

本条所列来源为《宋史》卷485〈夏国传〉；《辽史》《金史》相关本纪；黑水城西夏文文书与考古报告。它们能够支持特定的年序、制度语言、政治行动或局部遗存，却不能无条件化为社会总量。账本特别提醒“纪年框架可靠；该记录仅确证政权纪年连续，年内具体事项须由本年条另证”。因此，军报的兵数、条约的名分、官府的族群分类、判牍中的道德判断以及后出的叙述都须回到产生它们的场景，而不可直接代表全域动机或人口。

将事件与地方层次连接时，应同时考察资源怎样被征集、信息怎样传递、迁徙怎样改变户籍和劳动、宗教与亲属网络怎样提供或限制协作。现有证据往往只照亮少数地点；保留这种不均衡，才能避免把宋辽夏金元的多政权历史改写成单一中心的必然过程。

账本条目\texttt{XIA-0988-REGNAL-YEAR}所记988年的“党项李氏以李继迁在位第7年作为诏令、赋役与外交文书的纪年框架”，属于党项李氏的历史语境。条目所示前因是新岁颁历、年号和纪年是中枢与地方行政沟通的共同前提，其可见后果为为同年政令、战事和地方材料提供日期锚点，不能单独推知具体政策执行。这一变化若进入区域生活，必须通过道路、文书、军队、市场、家庭和地方中介才会发生；政权易手或法令发布不等于所有城镇、乡村与边地在同一时间获得同样经验。

本条所列来源为《宋史》卷485〈夏国传〉；《辽史》《金史》相关本纪；黑水城西夏文文书与考古报告。它们能够支持特定的年序、制度语言、政治行动或局部遗存，却不能无条件化为社会总量。账本特别提醒“纪年框架可靠；该记录仅确证政权纪年连续，年内具体事项须由本年条另证”。因此，军报的兵数、条约的名分、官府的族群分类、判牍中的道德判断以及后出的叙述都须回到产生它们的场景，而不可直接代表全域动机或人口。

将事件与地方层次连接时，应同时考察资源怎样被征集、信息怎样传递、迁徙怎样改变户籍和劳动、宗教与亲属网络怎样提供或限制协作。现有证据往往只照亮少数地点；保留这种不均衡，才能避免把宋辽夏金元的多政权历史改写成单一中心的必然过程。

账本条目\texttt{XIA-0989-REGNAL-YEAR}所记989年的“党项李氏以李继迁在位第8年作为诏令、赋役与外交文书的纪年框架”，属于党项李氏的历史语境。条目所示前因是新岁颁历、年号和纪年是中枢与地方行政沟通的共同前提，其可见后果为为同年政令、战事和地方材料提供日期锚点，不能单独推知具体政策执行。这一变化若进入区域生活，必须通过道路、文书、军队、市场、家庭和地方中介才会发生；政权易手或法令发布不等于所有城镇、乡村与边地在同一时间获得同样经验。

本条所列来源为《宋史》卷485〈夏国传〉；《辽史》《金史》相关本纪；黑水城西夏文文书与考古报告。它们能够支持特定的年序、制度语言、政治行动或局部遗存，却不能无条件化为社会总量。账本特别提醒“纪年框架可靠；该记录仅确证政权纪年连续，年内具体事项须由本年条另证”。因此，军报的兵数、条约的名分、官府的族群分类、判牍中的道德判断以及后出的叙述都须回到产生它们的场景，而不可直接代表全域动机或人口。

将事件与地方层次连接时，应同时考察资源怎样被征集、信息怎样传递、迁徙怎样改变户籍和劳动、宗教与亲属网络怎样提供或限制协作。现有证据往往只照亮少数地点；保留这种不均衡，才能避免把宋辽夏金元的多政权历史改写成单一中心的必然过程。

账本条目\texttt{XIA-0990-EVENT-050}所记990年的“李继迁在战与和之间接受宋廷封授并保持独立军事行动”，属于党项李氏与宋的历史语境。条目所示前因是前期边防、财政和统治联盟的压力构成直接背景，具体条件须按本条材料分辨，其可见后果为改变后续资源配置与政治选择；实际执行范围和社会影响仍须分项核验。这一变化若进入区域生活，必须通过道路、文书、军队、市场、家庭和地方中介才会发生；政权易手或法令发布不等于所有城镇、乡村与边地在同一时间获得同样经验。

本条所列来源为《宋史》卷485〈夏国传〉；《辽史》《金史》相关本纪；黑水城西夏文文书与考古报告。它们能够支持特定的年序、制度语言、政治行动或局部遗存，却不能无条件化为社会总量。账本特别提醒“编年主干可靠；细节、规模与动机须与对方材料、地方文书或考古材料互校”。因此，军报的兵数、条约的名分、官府的族群分类、判牍中的道德判断以及后出的叙述都须回到产生它们的场景，而不可直接代表全域动机或人口。

将事件与地方层次连接时，应同时考察资源怎样被征集、信息怎样传递、迁徙怎样改变户籍和劳动、宗教与亲属网络怎样提供或限制协作。现有证据往往只照亮少数地点；保留这种不均衡，才能避免把宋辽夏金元的多政权历史改写成单一中心的必然过程。

账本条目\texttt{XIA-0990-REGNAL-YEAR}所记990年的“党项李氏以李继迁在位第9年作为诏令、赋役与外交文书的纪年框架”，属于党项李氏的历史语境。条目所示前因是新岁颁历、年号和纪年是中枢与地方行政沟通的共同前提，其可见后果为为同年政令、战事和地方材料提供日期锚点，不能单独推知具体政策执行。这一变化若进入区域生活，必须通过道路、文书、军队、市场、家庭和地方中介才会发生；政权易手或法令发布不等于所有城镇、乡村与边地在同一时间获得同样经验。

本条所列来源为《宋史》卷485〈夏国传〉；《辽史》《金史》相关本纪；黑水城西夏文文书与考古报告。它们能够支持特定的年序、制度语言、政治行动或局部遗存，却不能无条件化为社会总量。账本特别提醒“纪年框架可靠；该记录仅确证政权纪年连续，年内具体事项须由本年条另证”。因此，军报的兵数、条约的名分、官府的族群分类、判牍中的道德判断以及后出的叙述都须回到产生它们的场景，而不可直接代表全域动机或人口。

将事件与地方层次连接时，应同时考察资源怎样被征集、信息怎样传递、迁徙怎样改变户籍和劳动、宗教与亲属网络怎样提供或限制协作。现有证据往往只照亮少数地点；保留这种不均衡，才能避免把宋辽夏金元的多政权历史改写成单一中心的必然过程。

账本条目\texttt{XIA-0991-REGNAL-YEAR}所记991年的“党项李氏以李继迁在位第10年作为诏令、赋役与外交文书的纪年框架”，属于党项李氏的历史语境。条目所示前因是新岁颁历、年号和纪年是中枢与地方行政沟通的共同前提，其可见后果为为同年政令、战事和地方材料提供日期锚点，不能单独推知具体政策执行。这一变化若进入区域生活，必须通过道路、文书、军队、市场、家庭和地方中介才会发生；政权易手或法令发布不等于所有城镇、乡村与边地在同一时间获得同样经验。

本条所列来源为《宋史》卷485〈夏国传〉；《辽史》《金史》相关本纪；黑水城西夏文文书与考古报告。它们能够支持特定的年序、制度语言、政治行动或局部遗存，却不能无条件化为社会总量。账本特别提醒“纪年框架可靠；该记录仅确证政权纪年连续，年内具体事项须由本年条另证”。因此，军报的兵数、条约的名分、官府的族群分类、判牍中的道德判断以及后出的叙述都须回到产生它们的场景，而不可直接代表全域动机或人口。

将事件与地方层次连接时，应同时考察资源怎样被征集、信息怎样传递、迁徙怎样改变户籍和劳动、宗教与亲属网络怎样提供或限制协作。现有证据往往只照亮少数地点；保留这种不均衡，才能避免把宋辽夏金元的多政权历史改写成单一中心的必然过程。

\subsection{外交、道路与人口移动}

账本条目\texttt{XIA-0992-REGNAL-YEAR}所记992年的“党项李氏以李继迁在位第11年作为诏令、赋役与外交文书的纪年框架”，属于党项李氏的历史语境。条目所示前因是新岁颁历、年号和纪年是中枢与地方行政沟通的共同前提，其可见后果为为同年政令、战事和地方材料提供日期锚点，不能单独推知具体政策执行。这一变化若进入区域生活，必须通过道路、文书、军队、市场、家庭和地方中介才会发生；政权易手或法令发布不等于所有城镇、乡村与边地在同一时间获得同样经验。

本条所列来源为《宋史》卷485〈夏国传〉；《辽史》《金史》相关本纪；黑水城西夏文文书与考古报告。它们能够支持特定的年序、制度语言、政治行动或局部遗存，却不能无条件化为社会总量。账本特别提醒“纪年框架可靠；该记录仅确证政权纪年连续，年内具体事项须由本年条另证”。因此，军报的兵数、条约的名分、官府的族群分类、判牍中的道德判断以及后出的叙述都须回到产生它们的场景，而不可直接代表全域动机或人口。

将事件与地方层次连接时，应同时考察资源怎样被征集、信息怎样传递、迁徙怎样改变户籍和劳动、宗教与亲属网络怎样提供或限制协作。现有证据往往只照亮少数地点；保留这种不均衡，才能避免把宋辽夏金元的多政权历史改写成单一中心的必然过程。

账本条目\texttt{XIA-0993-REGNAL-YEAR}所记993年的“党项李氏以李继迁在位第12年作为诏令、赋役与外交文书的纪年框架”，属于党项李氏的历史语境。条目所示前因是新岁颁历、年号和纪年是中枢与地方行政沟通的共同前提，其可见后果为为同年政令、战事和地方材料提供日期锚点，不能单独推知具体政策执行。这一变化若进入区域生活，必须通过道路、文书、军队、市场、家庭和地方中介才会发生；政权易手或法令发布不等于所有城镇、乡村与边地在同一时间获得同样经验。

本条所列来源为《宋史》卷485〈夏国传〉；《辽史》《金史》相关本纪；黑水城西夏文文书与考古报告。它们能够支持特定的年序、制度语言、政治行动或局部遗存，却不能无条件化为社会总量。账本特别提醒“纪年框架可靠；该记录仅确证政权纪年连续，年内具体事项须由本年条另证”。因此，军报的兵数、条约的名分、官府的族群分类、判牍中的道德判断以及后出的叙述都须回到产生它们的场景，而不可直接代表全域动机或人口。

将事件与地方层次连接时，应同时考察资源怎样被征集、信息怎样传递、迁徙怎样改变户籍和劳动、宗教与亲属网络怎样提供或限制协作。现有证据往往只照亮少数地点；保留这种不均衡，才能避免把宋辽夏金元的多政权历史改写成单一中心的必然过程。

账本条目\texttt{XIA-0994-REGNAL-YEAR}所记994年的“党项李氏以李继迁在位第13年作为诏令、赋役与外交文书的纪年框架”，属于党项李氏的历史语境。条目所示前因是新岁颁历、年号和纪年是中枢与地方行政沟通的共同前提，其可见后果为为同年政令、战事和地方材料提供日期锚点，不能单独推知具体政策执行。这一变化若进入区域生活，必须通过道路、文书、军队、市场、家庭和地方中介才会发生；政权易手或法令发布不等于所有城镇、乡村与边地在同一时间获得同样经验。

本条所列来源为《宋史》卷485〈夏国传〉；《辽史》《金史》相关本纪；黑水城西夏文文书与考古报告。它们能够支持特定的年序、制度语言、政治行动或局部遗存，却不能无条件化为社会总量。账本特别提醒“纪年框架可靠；该记录仅确证政权纪年连续，年内具体事项须由本年条另证”。因此，军报的兵数、条约的名分、官府的族群分类、判牍中的道德判断以及后出的叙述都须回到产生它们的场景，而不可直接代表全域动机或人口。

将事件与地方层次连接时，应同时考察资源怎样被征集、信息怎样传递、迁徙怎样改变户籍和劳动、宗教与亲属网络怎样提供或限制协作。现有证据往往只照亮少数地点；保留这种不均衡，才能避免把宋辽夏金元的多政权历史改写成单一中心的必然过程。

账本条目\texttt{XIA-0995-REGNAL-YEAR}所记995年的“党项李氏以李继迁在位第14年作为诏令、赋役与外交文书的纪年框架”，属于党项李氏的历史语境。条目所示前因是新岁颁历、年号和纪年是中枢与地方行政沟通的共同前提，其可见后果为为同年政令、战事和地方材料提供日期锚点，不能单独推知具体政策执行。这一变化若进入区域生活，必须通过道路、文书、军队、市场、家庭和地方中介才会发生；政权易手或法令发布不等于所有城镇、乡村与边地在同一时间获得同样经验。

本条所列来源为《宋史》卷485〈夏国传〉；《辽史》《金史》相关本纪；黑水城西夏文文书与考古报告。它们能够支持特定的年序、制度语言、政治行动或局部遗存，却不能无条件化为社会总量。账本特别提醒“纪年框架可靠；该记录仅确证政权纪年连续，年内具体事项须由本年条另证”。因此，军报的兵数、条约的名分、官府的族群分类、判牍中的道德判断以及后出的叙述都须回到产生它们的场景，而不可直接代表全域动机或人口。

将事件与地方层次连接时，应同时考察资源怎样被征集、信息怎样传递、迁徙怎样改变户籍和劳动、宗教与亲属网络怎样提供或限制协作。现有证据往往只照亮少数地点；保留这种不均衡，才能避免把宋辽夏金元的多政权历史改写成单一中心的必然过程。

账本条目\texttt{XIA-0996-REGNAL-YEAR}所记996年的“党项李氏以李继迁在位第15年作为诏令、赋役与外交文书的纪年框架”，属于党项李氏的历史语境。条目所示前因是新岁颁历、年号和纪年是中枢与地方行政沟通的共同前提，其可见后果为为同年政令、战事和地方材料提供日期锚点，不能单独推知具体政策执行。这一变化若进入区域生活，必须通过道路、文书、军队、市场、家庭和地方中介才会发生；政权易手或法令发布不等于所有城镇、乡村与边地在同一时间获得同样经验。

本条所列来源为《宋史》卷485〈夏国传〉；《辽史》《金史》相关本纪；黑水城西夏文文书与考古报告。它们能够支持特定的年序、制度语言、政治行动或局部遗存，却不能无条件化为社会总量。账本特别提醒“纪年框架可靠；该记录仅确证政权纪年连续，年内具体事项须由本年条另证”。因此，军报的兵数、条约的名分、官府的族群分类、判牍中的道德判断以及后出的叙述都须回到产生它们的场景，而不可直接代表全域动机或人口。

将事件与地方层次连接时，应同时考察资源怎样被征集、信息怎样传递、迁徙怎样改变户籍和劳动、宗教与亲属网络怎样提供或限制协作。现有证据往往只照亮少数地点；保留这种不均衡，才能避免把宋辽夏金元的多政权历史改写成单一中心的必然过程。

账本条目\texttt{XIA-0997-REGNAL-YEAR}所记997年的“党项李氏以李继迁在位第16年作为诏令、赋役与外交文书的纪年框架”，属于党项李氏的历史语境。条目所示前因是新岁颁历、年号和纪年是中枢与地方行政沟通的共同前提，其可见后果为为同年政令、战事和地方材料提供日期锚点，不能单独推知具体政策执行。这一变化若进入区域生活，必须通过道路、文书、军队、市场、家庭和地方中介才会发生；政权易手或法令发布不等于所有城镇、乡村与边地在同一时间获得同样经验。

本条所列来源为《宋史》卷485〈夏国传〉；《辽史》《金史》相关本纪；黑水城西夏文文书与考古报告。它们能够支持特定的年序、制度语言、政治行动或局部遗存，却不能无条件化为社会总量。账本特别提醒“纪年框架可靠；该记录仅确证政权纪年连续，年内具体事项须由本年条另证”。因此，军报的兵数、条约的名分、官府的族群分类、判牍中的道德判断以及后出的叙述都须回到产生它们的场景，而不可直接代表全域动机或人口。

将事件与地方层次连接时，应同时考察资源怎样被征集、信息怎样传递、迁徙怎样改变户籍和劳动、宗教与亲属网络怎样提供或限制协作。现有证据往往只照亮少数地点；保留这种不均衡，才能避免把宋辽夏金元的多政权历史改写成单一中心的必然过程。

账本条目\texttt{XIA-0998-REGNAL-YEAR}所记998年的“党项李氏以李继迁在位第17年作为诏令、赋役与外交文书的纪年框架”，属于党项李氏的历史语境。条目所示前因是新岁颁历、年号和纪年是中枢与地方行政沟通的共同前提，其可见后果为为同年政令、战事和地方材料提供日期锚点，不能单独推知具体政策执行。这一变化若进入区域生活，必须通过道路、文书、军队、市场、家庭和地方中介才会发生；政权易手或法令发布不等于所有城镇、乡村与边地在同一时间获得同样经验。

本条所列来源为《宋史》卷485〈夏国传〉；《辽史》《金史》相关本纪；黑水城西夏文文书与考古报告。它们能够支持特定的年序、制度语言、政治行动或局部遗存，却不能无条件化为社会总量。账本特别提醒“纪年框架可靠；该记录仅确证政权纪年连续，年内具体事项须由本年条另证”。因此，军报的兵数、条约的名分、官府的族群分类、判牍中的道德判断以及后出的叙述都须回到产生它们的场景，而不可直接代表全域动机或人口。

将事件与地方层次连接时，应同时考察资源怎样被征集、信息怎样传递、迁徙怎样改变户籍和劳动、宗教与亲属网络怎样提供或限制协作。现有证据往往只照亮少数地点；保留这种不均衡，才能避免把宋辽夏金元的多政权历史改写成单一中心的必然过程。

账本条目\texttt{XIA-0999-REGNAL-YEAR}所记999年的“党项李氏以李继迁在位第18年作为诏令、赋役与外交文书的纪年框架”，属于党项李氏的历史语境。条目所示前因是新岁颁历、年号和纪年是中枢与地方行政沟通的共同前提，其可见后果为为同年政令、战事和地方材料提供日期锚点，不能单独推知具体政策执行。这一变化若进入区域生活，必须通过道路、文书、军队、市场、家庭和地方中介才会发生；政权易手或法令发布不等于所有城镇、乡村与边地在同一时间获得同样经验。

本条所列来源为《宋史》卷485〈夏国传〉；《辽史》《金史》相关本纪；黑水城西夏文文书与考古报告。它们能够支持特定的年序、制度语言、政治行动或局部遗存，却不能无条件化为社会总量。账本特别提醒“纪年框架可靠；该记录仅确证政权纪年连续，年内具体事项须由本年条另证”。因此，军报的兵数、条约的名分、官府的族群分类、判牍中的道德判断以及后出的叙述都须回到产生它们的场景，而不可直接代表全域动机或人口。

将事件与地方层次连接时，应同时考察资源怎样被征集、信息怎样传递、迁徙怎样改变户籍和劳动、宗教与亲属网络怎样提供或限制协作。现有证据往往只照亮少数地点；保留这种不均衡，才能避免把宋辽夏金元的多政权历史改写成单一中心的必然过程。

账本条目\texttt{XIA-1000-REGNAL-YEAR}所记1000年的“党项李氏以李继迁在位第19年作为诏令、赋役与外交文书的纪年框架”，属于党项李氏的历史语境。条目所示前因是新岁颁历、年号和纪年是中枢与地方行政沟通的共同前提，其可见后果为为同年政令、战事和地方材料提供日期锚点，不能单独推知具体政策执行。这一变化若进入区域生活，必须通过道路、文书、军队、市场、家庭和地方中介才会发生；政权易手或法令发布不等于所有城镇、乡村与边地在同一时间获得同样经验。

本条所列来源为《宋史》卷485〈夏国传〉；《辽史》《金史》相关本纪；黑水城西夏文文书与考古报告。它们能够支持特定的年序、制度语言、政治行动或局部遗存，却不能无条件化为社会总量。账本特别提醒“纪年框架可靠；该记录仅确证政权纪年连续，年内具体事项须由本年条另证”。因此，军报的兵数、条约的名分、官府的族群分类、判牍中的道德判断以及后出的叙述都须回到产生它们的场景，而不可直接代表全域动机或人口。

将事件与地方层次连接时，应同时考察资源怎样被征集、信息怎样传递、迁徙怎样改变户籍和劳动、宗教与亲属网络怎样提供或限制协作。现有证据往往只照亮少数地点；保留这种不均衡，才能避免把宋辽夏金元的多政权历史改写成单一中心的必然过程。

账本条目\texttt{XIA-1001-REGNAL-YEAR}所记1001年的“党项李氏以李继迁在位第20年作为诏令、赋役与外交文书的纪年框架”，属于党项李氏的历史语境。条目所示前因是新岁颁历、年号和纪年是中枢与地方行政沟通的共同前提，其可见后果为为同年政令、战事和地方材料提供日期锚点，不能单独推知具体政策执行。这一变化若进入区域生活，必须通过道路、文书、军队、市场、家庭和地方中介才会发生；政权易手或法令发布不等于所有城镇、乡村与边地在同一时间获得同样经验。

本条所列来源为《宋史》卷485〈夏国传〉；《辽史》《金史》相关本纪；黑水城西夏文文书与考古报告。它们能够支持特定的年序、制度语言、政治行动或局部遗存，却不能无条件化为社会总量。账本特别提醒“纪年框架可靠；该记录仅确证政权纪年连续，年内具体事项须由本年条另证”。因此，军报的兵数、条约的名分、官府的族群分类、判牍中的道德判断以及后出的叙述都须回到产生它们的场景，而不可直接代表全域动机或人口。

将事件与地方层次连接时，应同时考察资源怎样被征集、信息怎样传递、迁徙怎样改变户籍和劳动、宗教与亲属网络怎样提供或限制协作。现有证据往往只照亮少数地点；保留这种不均衡，才能避免把宋辽夏金元的多政权历史改写成单一中心的必然过程。

账本条目\texttt{XIA-1002-REGNAL-YEAR}所记1002年的“党项李氏以李继迁在位第21年作为诏令、赋役与外交文书的纪年框架”，属于党项李氏的历史语境。条目所示前因是新岁颁历、年号和纪年是中枢与地方行政沟通的共同前提，其可见后果为为同年政令、战事和地方材料提供日期锚点，不能单独推知具体政策执行。这一变化若进入区域生活，必须通过道路、文书、军队、市场、家庭和地方中介才会发生；政权易手或法令发布不等于所有城镇、乡村与边地在同一时间获得同样经验。

本条所列来源为《宋史》卷485〈夏国传〉；《辽史》《金史》相关本纪；黑水城西夏文文书与考古报告。它们能够支持特定的年序、制度语言、政治行动或局部遗存，却不能无条件化为社会总量。账本特别提醒“纪年框架可靠；该记录仅确证政权纪年连续，年内具体事项须由本年条另证”。因此，军报的兵数、条约的名分、官府的族群分类、判牍中的道德判断以及后出的叙述都须回到产生它们的场景，而不可直接代表全域动机或人口。

将事件与地方层次连接时，应同时考察资源怎样被征集、信息怎样传递、迁徙怎样改变户籍和劳动、宗教与亲属网络怎样提供或限制协作。现有证据往往只照亮少数地点；保留这种不均衡，才能避免把宋辽夏金元的多政权历史改写成单一中心的必然过程。

账本条目\texttt{XIA-1003-REGNAL-YEAR}所记1003年的“党项李氏以李继迁在位第22年作为诏令、赋役与外交文书的纪年框架”，属于党项李氏的历史语境。条目所示前因是新岁颁历、年号和纪年是中枢与地方行政沟通的共同前提，其可见后果为为同年政令、战事和地方材料提供日期锚点，不能单独推知具体政策执行。这一变化若进入区域生活，必须通过道路、文书、军队、市场、家庭和地方中介才会发生；政权易手或法令发布不等于所有城镇、乡村与边地在同一时间获得同样经验。

本条所列来源为《宋史》卷485〈夏国传〉；《辽史》《金史》相关本纪；黑水城西夏文文书与考古报告。它们能够支持特定的年序、制度语言、政治行动或局部遗存，却不能无条件化为社会总量。账本特别提醒“纪年框架可靠；该记录仅确证政权纪年连续，年内具体事项须由本年条另证”。因此，军报的兵数、条约的名分、官府的族群分类、判牍中的道德判断以及后出的叙述都须回到产生它们的场景，而不可直接代表全域动机或人口。

将事件与地方层次连接时，应同时考察资源怎样被征集、信息怎样传递、迁徙怎样改变户籍和劳动、宗教与亲属网络怎样提供或限制协作。现有证据往往只照亮少数地点；保留这种不均衡，才能避免把宋辽夏金元的多政权历史改写成单一中心的必然过程。

\subsection{环境风险与地方经济}

账本条目\texttt{XIA-1004-EVENT-051}所记1004年的“李继迁去世，李德明继承部众”，属于党项李氏的历史语境。条目所示前因是前期边防、财政和统治联盟的压力构成直接背景，具体条件须按本条材料分辨，其可见后果为改变后续资源配置与政治选择；实际执行范围和社会影响仍须分项核验。这一变化若进入区域生活，必须通过道路、文书、军队、市场、家庭和地方中介才会发生；政权易手或法令发布不等于所有城镇、乡村与边地在同一时间获得同样经验。

本条所列来源为《宋史》卷485〈夏国传〉；《辽史》《金史》相关本纪；黑水城西夏文文书与考古报告。它们能够支持特定的年序、制度语言、政治行动或局部遗存，却不能无条件化为社会总量。账本特别提醒“编年主干可靠；细节、规模与动机须与对方材料、地方文书或考古材料互校”。因此，军报的兵数、条约的名分、官府的族群分类、判牍中的道德判断以及后出的叙述都须回到产生它们的场景，而不可直接代表全域动机或人口。

将事件与地方层次连接时，应同时考察资源怎样被征集、信息怎样传递、迁徙怎样改变户籍和劳动、宗教与亲属网络怎样提供或限制协作。现有证据往往只照亮少数地点；保留这种不均衡，才能避免把宋辽夏金元的多政权历史改写成单一中心的必然过程。

账本条目\texttt{XIA-1004-REGNAL-YEAR}所记1004年的“党项李氏以李继迁在位第23年作为诏令、赋役与外交文书的纪年框架”，属于党项李氏的历史语境。条目所示前因是新岁颁历、年号和纪年是中枢与地方行政沟通的共同前提，其可见后果为为同年政令、战事和地方材料提供日期锚点，不能单独推知具体政策执行。这一变化若进入区域生活，必须通过道路、文书、军队、市场、家庭和地方中介才会发生；政权易手或法令发布不等于所有城镇、乡村与边地在同一时间获得同样经验。

本条所列来源为《宋史》卷485〈夏国传〉；《辽史》《金史》相关本纪；黑水城西夏文文书与考古报告。它们能够支持特定的年序、制度语言、政治行动或局部遗存，却不能无条件化为社会总量。账本特别提醒“纪年框架可靠；该记录仅确证政权纪年连续，年内具体事项须由本年条另证”。因此，军报的兵数、条约的名分、官府的族群分类、判牍中的道德判断以及后出的叙述都须回到产生它们的场景，而不可直接代表全域动机或人口。

将事件与地方层次连接时，应同时考察资源怎样被征集、信息怎样传递、迁徙怎样改变户籍和劳动、宗教与亲属网络怎样提供或限制协作。现有证据往往只照亮少数地点；保留这种不均衡，才能避免把宋辽夏金元的多政权历史改写成单一中心的必然过程。

账本条目\texttt{XIA-1005-REGNAL-YEAR}所记1005年的“党项李氏以李德明在位第1年作为诏令、赋役与外交文书的纪年框架”，属于党项李氏的历史语境。条目所示前因是新岁颁历、年号和纪年是中枢与地方行政沟通的共同前提，其可见后果为为同年政令、战事和地方材料提供日期锚点，不能单独推知具体政策执行。这一变化若进入区域生活，必须通过道路、文书、军队、市场、家庭和地方中介才会发生；政权易手或法令发布不等于所有城镇、乡村与边地在同一时间获得同样经验。

本条所列来源为《宋史》卷485〈夏国传〉；《辽史》《金史》相关本纪；黑水城西夏文文书与考古报告。它们能够支持特定的年序、制度语言、政治行动或局部遗存，却不能无条件化为社会总量。账本特别提醒“纪年框架可靠；该记录仅确证政权纪年连续，年内具体事项须由本年条另证”。因此，军报的兵数、条约的名分、官府的族群分类、判牍中的道德判断以及后出的叙述都须回到产生它们的场景，而不可直接代表全域动机或人口。

将事件与地方层次连接时，应同时考察资源怎样被征集、信息怎样传递、迁徙怎样改变户籍和劳动、宗教与亲属网络怎样提供或限制协作。现有证据往往只照亮少数地点；保留这种不均衡，才能避免把宋辽夏金元的多政权历史改写成单一中心的必然过程。

账本条目\texttt{XIA-1006-REGNAL-YEAR}所记1006年的“党项李氏以李德明在位第2年作为诏令、赋役与外交文书的纪年框架”，属于党项李氏的历史语境。条目所示前因是新岁颁历、年号和纪年是中枢与地方行政沟通的共同前提，其可见后果为为同年政令、战事和地方材料提供日期锚点，不能单独推知具体政策执行。这一变化若进入区域生活，必须通过道路、文书、军队、市场、家庭和地方中介才会发生；政权易手或法令发布不等于所有城镇、乡村与边地在同一时间获得同样经验。

本条所列来源为《宋史》卷485〈夏国传〉；《辽史》《金史》相关本纪；黑水城西夏文文书与考古报告。它们能够支持特定的年序、制度语言、政治行动或局部遗存，却不能无条件化为社会总量。账本特别提醒“纪年框架可靠；该记录仅确证政权纪年连续，年内具体事项须由本年条另证”。因此，军报的兵数、条约的名分、官府的族群分类、判牍中的道德判断以及后出的叙述都须回到产生它们的场景，而不可直接代表全域动机或人口。

将事件与地方层次连接时，应同时考察资源怎样被征集、信息怎样传递、迁徙怎样改变户籍和劳动、宗教与亲属网络怎样提供或限制协作。现有证据往往只照亮少数地点；保留这种不均衡，才能避免把宋辽夏金元的多政权历史改写成单一中心的必然过程。

账本条目\texttt{XIA-1007-REGNAL-YEAR}所记1007年的“党项李氏以李德明在位第3年作为诏令、赋役与外交文书的纪年框架”，属于党项李氏的历史语境。条目所示前因是新岁颁历、年号和纪年是中枢与地方行政沟通的共同前提，其可见后果为为同年政令、战事和地方材料提供日期锚点，不能单独推知具体政策执行。这一变化若进入区域生活，必须通过道路、文书、军队、市场、家庭和地方中介才会发生；政权易手或法令发布不等于所有城镇、乡村与边地在同一时间获得同样经验。

本条所列来源为《宋史》卷485〈夏国传〉；《辽史》《金史》相关本纪；黑水城西夏文文书与考古报告。它们能够支持特定的年序、制度语言、政治行动或局部遗存，却不能无条件化为社会总量。账本特别提醒“纪年框架可靠；该记录仅确证政权纪年连续，年内具体事项须由本年条另证”。因此，军报的兵数、条约的名分、官府的族群分类、判牍中的道德判断以及后出的叙述都须回到产生它们的场景，而不可直接代表全域动机或人口。

将事件与地方层次连接时，应同时考察资源怎样被征集、信息怎样传递、迁徙怎样改变户籍和劳动、宗教与亲属网络怎样提供或限制协作。现有证据往往只照亮少数地点；保留这种不均衡，才能避免把宋辽夏金元的多政权历史改写成单一中心的必然过程。

账本条目\texttt{XIA-1008-REGNAL-YEAR}所记1008年的“党项李氏以李德明在位第4年作为诏令、赋役与外交文书的纪年框架”，属于党项李氏的历史语境。条目所示前因是新岁颁历、年号和纪年是中枢与地方行政沟通的共同前提，其可见后果为为同年政令、战事和地方材料提供日期锚点，不能单独推知具体政策执行。这一变化若进入区域生活，必须通过道路、文书、军队、市场、家庭和地方中介才会发生；政权易手或法令发布不等于所有城镇、乡村与边地在同一时间获得同样经验。

本条所列来源为《宋史》卷485〈夏国传〉；《辽史》《金史》相关本纪；黑水城西夏文文书与考古报告。它们能够支持特定的年序、制度语言、政治行动或局部遗存，却不能无条件化为社会总量。账本特别提醒“纪年框架可靠；该记录仅确证政权纪年连续，年内具体事项须由本年条另证”。因此，军报的兵数、条约的名分、官府的族群分类、判牍中的道德判断以及后出的叙述都须回到产生它们的场景，而不可直接代表全域动机或人口。

将事件与地方层次连接时，应同时考察资源怎样被征集、信息怎样传递、迁徙怎样改变户籍和劳动、宗教与亲属网络怎样提供或限制协作。现有证据往往只照亮少数地点；保留这种不均衡，才能避免把宋辽夏金元的多政权历史改写成单一中心的必然过程。

账本条目\texttt{XIA-1009-EVENT-052}所记1009年的“李德明接受宋廷名号并经营夏州、灵州一带”，属于党项李氏与宋的历史语境。条目所示前因是前期边防、财政和统治联盟的压力构成直接背景，具体条件须按本条材料分辨，其可见后果为改变后续资源配置与政治选择；实际执行范围和社会影响仍须分项核验。这一变化若进入区域生活，必须通过道路、文书、军队、市场、家庭和地方中介才会发生；政权易手或法令发布不等于所有城镇、乡村与边地在同一时间获得同样经验。

本条所列来源为《宋史》卷485〈夏国传〉；《辽史》《金史》相关本纪；黑水城西夏文文书与考古报告。它们能够支持特定的年序、制度语言、政治行动或局部遗存，却不能无条件化为社会总量。账本特别提醒“编年主干可靠；细节、规模与动机须与对方材料、地方文书或考古材料互校”。因此，军报的兵数、条约的名分、官府的族群分类、判牍中的道德判断以及后出的叙述都须回到产生它们的场景，而不可直接代表全域动机或人口。

将事件与地方层次连接时，应同时考察资源怎样被征集、信息怎样传递、迁徙怎样改变户籍和劳动、宗教与亲属网络怎样提供或限制协作。现有证据往往只照亮少数地点；保留这种不均衡，才能避免把宋辽夏金元的多政权历史改写成单一中心的必然过程。

账本条目\texttt{XIA-1009-REGNAL-YEAR}所记1009年的“党项李氏以李德明在位第5年作为诏令、赋役与外交文书的纪年框架”，属于党项李氏的历史语境。条目所示前因是新岁颁历、年号和纪年是中枢与地方行政沟通的共同前提，其可见后果为为同年政令、战事和地方材料提供日期锚点，不能单独推知具体政策执行。这一变化若进入区域生活，必须通过道路、文书、军队、市场、家庭和地方中介才会发生；政权易手或法令发布不等于所有城镇、乡村与边地在同一时间获得同样经验。

本条所列来源为《宋史》卷485〈夏国传〉；《辽史》《金史》相关本纪；黑水城西夏文文书与考古报告。它们能够支持特定的年序、制度语言、政治行动或局部遗存，却不能无条件化为社会总量。账本特别提醒“纪年框架可靠；该记录仅确证政权纪年连续，年内具体事项须由本年条另证”。因此，军报的兵数、条约的名分、官府的族群分类、判牍中的道德判断以及后出的叙述都须回到产生它们的场景，而不可直接代表全域动机或人口。

将事件与地方层次连接时，应同时考察资源怎样被征集、信息怎样传递、迁徙怎样改变户籍和劳动、宗教与亲属网络怎样提供或限制协作。现有证据往往只照亮少数地点；保留这种不均衡，才能避免把宋辽夏金元的多政权历史改写成单一中心的必然过程。

账本条目\texttt{XIA-1010-REGNAL-YEAR}所记1010年的“党项李氏以李德明在位第6年作为诏令、赋役与外交文书的纪年框架”，属于党项李氏的历史语境。条目所示前因是新岁颁历、年号和纪年是中枢与地方行政沟通的共同前提，其可见后果为为同年政令、战事和地方材料提供日期锚点，不能单独推知具体政策执行。这一变化若进入区域生活，必须通过道路、文书、军队、市场、家庭和地方中介才会发生；政权易手或法令发布不等于所有城镇、乡村与边地在同一时间获得同样经验。

本条所列来源为《宋史》卷485〈夏国传〉；《辽史》《金史》相关本纪；黑水城西夏文文书与考古报告。它们能够支持特定的年序、制度语言、政治行动或局部遗存，却不能无条件化为社会总量。账本特别提醒“纪年框架可靠；该记录仅确证政权纪年连续，年内具体事项须由本年条另证”。因此，军报的兵数、条约的名分、官府的族群分类、判牍中的道德判断以及后出的叙述都须回到产生它们的场景，而不可直接代表全域动机或人口。

将事件与地方层次连接时，应同时考察资源怎样被征集、信息怎样传递、迁徙怎样改变户籍和劳动、宗教与亲属网络怎样提供或限制协作。现有证据往往只照亮少数地点；保留这种不均衡，才能避免把宋辽夏金元的多政权历史改写成单一中心的必然过程。

账本条目\texttt{XIA-1011-REGNAL-YEAR}所记1011年的“党项李氏以李德明在位第7年作为诏令、赋役与外交文书的纪年框架”，属于党项李氏的历史语境。条目所示前因是新岁颁历、年号和纪年是中枢与地方行政沟通的共同前提，其可见后果为为同年政令、战事和地方材料提供日期锚点，不能单独推知具体政策执行。这一变化若进入区域生活，必须通过道路、文书、军队、市场、家庭和地方中介才会发生；政权易手或法令发布不等于所有城镇、乡村与边地在同一时间获得同样经验。

本条所列来源为《宋史》卷485〈夏国传〉；《辽史》《金史》相关本纪；黑水城西夏文文书与考古报告。它们能够支持特定的年序、制度语言、政治行动或局部遗存，却不能无条件化为社会总量。账本特别提醒“纪年框架可靠；该记录仅确证政权纪年连续，年内具体事项须由本年条另证”。因此，军报的兵数、条约的名分、官府的族群分类、判牍中的道德判断以及后出的叙述都须回到产生它们的场景，而不可直接代表全域动机或人口。

将事件与地方层次连接时，应同时考察资源怎样被征集、信息怎样传递、迁徙怎样改变户籍和劳动、宗教与亲属网络怎样提供或限制协作。现有证据往往只照亮少数地点；保留这种不均衡，才能避免把宋辽夏金元的多政权历史改写成单一中心的必然过程。

账本条目\texttt{XIA-1012-REGNAL-YEAR}所记1012年的“党项李氏以李德明在位第8年作为诏令、赋役与外交文书的纪年框架”，属于党项李氏的历史语境。条目所示前因是新岁颁历、年号和纪年是中枢与地方行政沟通的共同前提，其可见后果为为同年政令、战事和地方材料提供日期锚点，不能单独推知具体政策执行。这一变化若进入区域生活，必须通过道路、文书、军队、市场、家庭和地方中介才会发生；政权易手或法令发布不等于所有城镇、乡村与边地在同一时间获得同样经验。

本条所列来源为《宋史》卷485〈夏国传〉；《辽史》《金史》相关本纪；黑水城西夏文文书与考古报告。它们能够支持特定的年序、制度语言、政治行动或局部遗存，却不能无条件化为社会总量。账本特别提醒“纪年框架可靠；该记录仅确证政权纪年连续，年内具体事项须由本年条另证”。因此，军报的兵数、条约的名分、官府的族群分类、判牍中的道德判断以及后出的叙述都须回到产生它们的场景，而不可直接代表全域动机或人口。

将事件与地方层次连接时，应同时考察资源怎样被征集、信息怎样传递、迁徙怎样改变户籍和劳动、宗教与亲属网络怎样提供或限制协作。现有证据往往只照亮少数地点；保留这种不均衡，才能避免把宋辽夏金元的多政权历史改写成单一中心的必然过程。

账本条目\texttt{XIA-1013-REGNAL-YEAR}所记1013年的“党项李氏以李德明在位第9年作为诏令、赋役与外交文书的纪年框架”，属于党项李氏的历史语境。条目所示前因是新岁颁历、年号和纪年是中枢与地方行政沟通的共同前提，其可见后果为为同年政令、战事和地方材料提供日期锚点，不能单独推知具体政策执行。这一变化若进入区域生活，必须通过道路、文书、军队、市场、家庭和地方中介才会发生；政权易手或法令发布不等于所有城镇、乡村与边地在同一时间获得同样经验。

本条所列来源为《宋史》卷485〈夏国传〉；《辽史》《金史》相关本纪；黑水城西夏文文书与考古报告。它们能够支持特定的年序、制度语言、政治行动或局部遗存，却不能无条件化为社会总量。账本特别提醒“纪年框架可靠；该记录仅确证政权纪年连续，年内具体事项须由本年条另证”。因此，军报的兵数、条约的名分、官府的族群分类、判牍中的道德判断以及后出的叙述都须回到产生它们的场景，而不可直接代表全域动机或人口。

将事件与地方层次连接时，应同时考察资源怎样被征集、信息怎样传递、迁徙怎样改变户籍和劳动、宗教与亲属网络怎样提供或限制协作。现有证据往往只照亮少数地点；保留这种不均衡，才能避免把宋辽夏金元的多政权历史改写成单一中心的必然过程。

\subsection{宗教、法律与材料边界}

账本条目\texttt{XIA-1014-REGNAL-YEAR}所记1014年的“党项李氏以李德明在位第10年作为诏令、赋役与外交文书的纪年框架”，属于党项李氏的历史语境。条目所示前因是新岁颁历、年号和纪年是中枢与地方行政沟通的共同前提，其可见后果为为同年政令、战事和地方材料提供日期锚点，不能单独推知具体政策执行。这一变化若进入区域生活，必须通过道路、文书、军队、市场、家庭和地方中介才会发生；政权易手或法令发布不等于所有城镇、乡村与边地在同一时间获得同样经验。

本条所列来源为《宋史》卷485〈夏国传〉；《辽史》《金史》相关本纪；黑水城西夏文文书与考古报告。它们能够支持特定的年序、制度语言、政治行动或局部遗存，却不能无条件化为社会总量。账本特别提醒“纪年框架可靠；该记录仅确证政权纪年连续，年内具体事项须由本年条另证”。因此，军报的兵数、条约的名分、官府的族群分类、判牍中的道德判断以及后出的叙述都须回到产生它们的场景，而不可直接代表全域动机或人口。

将事件与地方层次连接时，应同时考察资源怎样被征集、信息怎样传递、迁徙怎样改变户籍和劳动、宗教与亲属网络怎样提供或限制协作。现有证据往往只照亮少数地点；保留这种不均衡，才能避免把宋辽夏金元的多政权历史改写成单一中心的必然过程。

账本条目\texttt{XIA-1015-REGNAL-YEAR}所记1015年的“党项李氏以李德明在位第11年作为诏令、赋役与外交文书的纪年框架”，属于党项李氏的历史语境。条目所示前因是新岁颁历、年号和纪年是中枢与地方行政沟通的共同前提，其可见后果为为同年政令、战事和地方材料提供日期锚点，不能单独推知具体政策执行。这一变化若进入区域生活，必须通过道路、文书、军队、市场、家庭和地方中介才会发生；政权易手或法令发布不等于所有城镇、乡村与边地在同一时间获得同样经验。

本条所列来源为《宋史》卷485〈夏国传〉；《辽史》《金史》相关本纪；黑水城西夏文文书与考古报告。它们能够支持特定的年序、制度语言、政治行动或局部遗存，却不能无条件化为社会总量。账本特别提醒“纪年框架可靠；该记录仅确证政权纪年连续，年内具体事项须由本年条另证”。因此，军报的兵数、条约的名分、官府的族群分类、判牍中的道德判断以及后出的叙述都须回到产生它们的场景，而不可直接代表全域动机或人口。

将事件与地方层次连接时，应同时考察资源怎样被征集、信息怎样传递、迁徙怎样改变户籍和劳动、宗教与亲属网络怎样提供或限制协作。现有证据往往只照亮少数地点；保留这种不均衡，才能避免把宋辽夏金元的多政权历史改写成单一中心的必然过程。

账本条目\texttt{XIA-1016-REGNAL-YEAR}所记1016年的“党项李氏以李德明在位第12年作为诏令、赋役与外交文书的纪年框架”，属于党项李氏的历史语境。条目所示前因是新岁颁历、年号和纪年是中枢与地方行政沟通的共同前提，其可见后果为为同年政令、战事和地方材料提供日期锚点，不能单独推知具体政策执行。这一变化若进入区域生活，必须通过道路、文书、军队、市场、家庭和地方中介才会发生；政权易手或法令发布不等于所有城镇、乡村与边地在同一时间获得同样经验。

本条所列来源为《宋史》卷485〈夏国传〉；《辽史》《金史》相关本纪；黑水城西夏文文书与考古报告。它们能够支持特定的年序、制度语言、政治行动或局部遗存，却不能无条件化为社会总量。账本特别提醒“纪年框架可靠；该记录仅确证政权纪年连续，年内具体事项须由本年条另证”。因此，军报的兵数、条约的名分、官府的族群分类、判牍中的道德判断以及后出的叙述都须回到产生它们的场景，而不可直接代表全域动机或人口。

将事件与地方层次连接时，应同时考察资源怎样被征集、信息怎样传递、迁徙怎样改变户籍和劳动、宗教与亲属网络怎样提供或限制协作。现有证据往往只照亮少数地点；保留这种不均衡，才能避免把宋辽夏金元的多政权历史改写成单一中心的必然过程。

账本条目\texttt{XIA-1017-REGNAL-YEAR}所记1017年的“党项李氏以李德明在位第13年作为诏令、赋役与外交文书的纪年框架”，属于党项李氏的历史语境。条目所示前因是新岁颁历、年号和纪年是中枢与地方行政沟通的共同前提，其可见后果为为同年政令、战事和地方材料提供日期锚点，不能单独推知具体政策执行。这一变化若进入区域生活，必须通过道路、文书、军队、市场、家庭和地方中介才会发生；政权易手或法令发布不等于所有城镇、乡村与边地在同一时间获得同样经验。

本条所列来源为《宋史》卷485〈夏国传〉；《辽史》《金史》相关本纪；黑水城西夏文文书与考古报告。它们能够支持特定的年序、制度语言、政治行动或局部遗存，却不能无条件化为社会总量。账本特别提醒“纪年框架可靠；该记录仅确证政权纪年连续，年内具体事项须由本年条另证”。因此，军报的兵数、条约的名分、官府的族群分类、判牍中的道德判断以及后出的叙述都须回到产生它们的场景，而不可直接代表全域动机或人口。

将事件与地方层次连接时，应同时考察资源怎样被征集、信息怎样传递、迁徙怎样改变户籍和劳动、宗教与亲属网络怎样提供或限制协作。现有证据往往只照亮少数地点；保留这种不均衡，才能避免把宋辽夏金元的多政权历史改写成单一中心的必然过程。

账本条目\texttt{XIA-1018-REGNAL-YEAR}所记1018年的“党项李氏以李德明在位第14年作为诏令、赋役与外交文书的纪年框架”，属于党项李氏的历史语境。条目所示前因是新岁颁历、年号和纪年是中枢与地方行政沟通的共同前提，其可见后果为为同年政令、战事和地方材料提供日期锚点，不能单独推知具体政策执行。这一变化若进入区域生活，必须通过道路、文书、军队、市场、家庭和地方中介才会发生；政权易手或法令发布不等于所有城镇、乡村与边地在同一时间获得同样经验。

本条所列来源为《宋史》卷485〈夏国传〉；《辽史》《金史》相关本纪；黑水城西夏文文书与考古报告。它们能够支持特定的年序、制度语言、政治行动或局部遗存，却不能无条件化为社会总量。账本特别提醒“纪年框架可靠；该记录仅确证政权纪年连续，年内具体事项须由本年条另证”。因此，军报的兵数、条约的名分、官府的族群分类、判牍中的道德判断以及后出的叙述都须回到产生它们的场景，而不可直接代表全域动机或人口。

将事件与地方层次连接时，应同时考察资源怎样被征集、信息怎样传递、迁徙怎样改变户籍和劳动、宗教与亲属网络怎样提供或限制协作。现有证据往往只照亮少数地点；保留这种不均衡，才能避免把宋辽夏金元的多政权历史改写成单一中心的必然过程。

账本条目\texttt{XIA-1019-REGNAL-YEAR}所记1019年的“党项李氏以李德明在位第15年作为诏令、赋役与外交文书的纪年框架”，属于党项李氏的历史语境。条目所示前因是新岁颁历、年号和纪年是中枢与地方行政沟通的共同前提，其可见后果为为同年政令、战事和地方材料提供日期锚点，不能单独推知具体政策执行。这一变化若进入区域生活，必须通过道路、文书、军队、市场、家庭和地方中介才会发生；政权易手或法令发布不等于所有城镇、乡村与边地在同一时间获得同样经验。

本条所列来源为《宋史》卷485〈夏国传〉；《辽史》《金史》相关本纪；黑水城西夏文文书与考古报告。它们能够支持特定的年序、制度语言、政治行动或局部遗存，却不能无条件化为社会总量。账本特别提醒“纪年框架可靠；该记录仅确证政权纪年连续，年内具体事项须由本年条另证”。因此，军报的兵数、条约的名分、官府的族群分类、判牍中的道德判断以及后出的叙述都须回到产生它们的场景，而不可直接代表全域动机或人口。

将事件与地方层次连接时，应同时考察资源怎样被征集、信息怎样传递、迁徙怎样改变户籍和劳动、宗教与亲属网络怎样提供或限制协作。现有证据往往只照亮少数地点；保留这种不均衡，才能避免把宋辽夏金元的多政权历史改写成单一中心的必然过程。

账本条目\texttt{XIA-1020-REGNAL-YEAR}所记1020年的“党项李氏以李德明在位第16年作为诏令、赋役与外交文书的纪年框架”，属于党项李氏的历史语境。条目所示前因是新岁颁历、年号和纪年是中枢与地方行政沟通的共同前提，其可见后果为为同年政令、战事和地方材料提供日期锚点，不能单独推知具体政策执行。这一变化若进入区域生活，必须通过道路、文书、军队、市场、家庭和地方中介才会发生；政权易手或法令发布不等于所有城镇、乡村与边地在同一时间获得同样经验。

本条所列来源为《宋史》卷485〈夏国传〉；《辽史》《金史》相关本纪；黑水城西夏文文书与考古报告。它们能够支持特定的年序、制度语言、政治行动或局部遗存，却不能无条件化为社会总量。账本特别提醒“纪年框架可靠；该记录仅确证政权纪年连续，年内具体事项须由本年条另证”。因此，军报的兵数、条约的名分、官府的族群分类、判牍中的道德判断以及后出的叙述都须回到产生它们的场景，而不可直接代表全域动机或人口。

将事件与地方层次连接时，应同时考察资源怎样被征集、信息怎样传递、迁徙怎样改变户籍和劳动、宗教与亲属网络怎样提供或限制协作。现有证据往往只照亮少数地点；保留这种不均衡，才能避免把宋辽夏金元的多政权历史改写成单一中心的必然过程。

账本条目\texttt{XIA-1021-REGNAL-YEAR}所记1021年的“党项李氏以李德明在位第17年作为诏令、赋役与外交文书的纪年框架”，属于党项李氏的历史语境。条目所示前因是新岁颁历、年号和纪年是中枢与地方行政沟通的共同前提，其可见后果为为同年政令、战事和地方材料提供日期锚点，不能单独推知具体政策执行。这一变化若进入区域生活，必须通过道路、文书、军队、市场、家庭和地方中介才会发生；政权易手或法令发布不等于所有城镇、乡村与边地在同一时间获得同样经验。

本条所列来源为《宋史》卷485〈夏国传〉；《辽史》《金史》相关本纪；黑水城西夏文文书与考古报告。它们能够支持特定的年序、制度语言、政治行动或局部遗存，却不能无条件化为社会总量。账本特别提醒“纪年框架可靠；该记录仅确证政权纪年连续，年内具体事项须由本年条另证”。因此，军报的兵数、条约的名分、官府的族群分类、判牍中的道德判断以及后出的叙述都须回到产生它们的场景，而不可直接代表全域动机或人口。

将事件与地方层次连接时，应同时考察资源怎样被征集、信息怎样传递、迁徙怎样改变户籍和劳动、宗教与亲属网络怎样提供或限制协作。现有证据往往只照亮少数地点；保留这种不均衡，才能避免把宋辽夏金元的多政权历史改写成单一中心的必然过程。

账本条目\texttt{XIA-1022-REGNAL-YEAR}所记1022年的“党项李氏以李德明在位第18年作为诏令、赋役与外交文书的纪年框架”，属于党项李氏的历史语境。条目所示前因是新岁颁历、年号和纪年是中枢与地方行政沟通的共同前提，其可见后果为为同年政令、战事和地方材料提供日期锚点，不能单独推知具体政策执行。这一变化若进入区域生活，必须通过道路、文书、军队、市场、家庭和地方中介才会发生；政权易手或法令发布不等于所有城镇、乡村与边地在同一时间获得同样经验。

本条所列来源为《宋史》卷485〈夏国传〉；《辽史》《金史》相关本纪；黑水城西夏文文书与考古报告。它们能够支持特定的年序、制度语言、政治行动或局部遗存，却不能无条件化为社会总量。账本特别提醒“纪年框架可靠；该记录仅确证政权纪年连续，年内具体事项须由本年条另证”。因此，军报的兵数、条约的名分、官府的族群分类、判牍中的道德判断以及后出的叙述都须回到产生它们的场景，而不可直接代表全域动机或人口。

将事件与地方层次连接时，应同时考察资源怎样被征集、信息怎样传递、迁徙怎样改变户籍和劳动、宗教与亲属网络怎样提供或限制协作。现有证据往往只照亮少数地点；保留这种不均衡，才能避免把宋辽夏金元的多政权历史改写成单一中心的必然过程。
